\documentclass[12pt]{article}

% ---------------------------------------------------------------
% Packages
% ---------------------------------------------------------------
\usepackage[margin=1in]{geometry}
\usepackage{setspace}
\usepackage{amsmath, amssymb}
\usepackage{graphicx}
\usepackage{booktabs}
\usepackage{threeparttable}
\usepackage{caption}
\usepackage{subcaption}
\usepackage{hyperref}
\usepackage{natbib}
\usepackage{placeins}
\usepackage{authblk}

\setcitestyle{authoryear, open={(}, close={)}}
\onehalfspacing

% ---------------------------------------------------------------
% Title / Authors
% ---------------------------------------------------------------

\title{Reading the Storm: Media Attribution of United States Climate Disasters, A Long Panel}

\author[1]{Hariaksha Gunda}
\author[2]{Dr. Michael Price}
\author[3]{Dr. Angela Doku}
\author[4]{Dr. John List}

\affil[1]{University of Alabama}
\affil[2]{University of Alabama}
\affil[3]{University of Toronto}
\affil[4]{University of Chicago}

\date{\today}

\begin{document}

\maketitle

% ---------------------------------------------------------------
\begin{abstract}
\begin{sloppypar}
Does the media tell readers why a disaster happened? We study how
U.S.\ newspapers frame the causal origins of billion-dollar climate
disasters, constructing an index of attribution that ranges from
explicit denial of a climate link to explicit anthropogenic
attribution, applied to articles covering every major disaster the
federal government has recorded since 2000. We document three
findings. First, silence is the norm: three in four articles make no
causal claim at all, and the typical disaster is covered in terms
close to causally neutral. Second, this silence is receding --
attribution framing has risen steadily over the past two decades,
consistent with a broader partisan realignment around climate change
that took hold after 2010. Third, an outlet's political lean strongly
predicts how it frames disaster causation, with the effect concentrated
in national and wire-service coverage rather than local newspapers.
By contrast, the severity, type, and location of a disaster are
essentially uninformative about how it gets framed. These patterns
suggest that public understanding of climate risk is shaped less by
the disasters themselves than by the media environment interpreting
them.
\end{sloppypar}
\end{abstract}

\noindent \textbf{Keywords:} climate change attribution, media framing,
natural language processing, disaster economics

\newpage


% ===================================================================
\section{Introduction}
% ===================================================================

Between 2000 and 2024, the United States experienced 313 weather and
climate disasters each causing at least one billion dollars in
CPI-adjusted losses, together inflicting more than \$2.36 trillion in
damages and 10,849 fatalities.  The pace has been accelerating:
\citet{boustan2020} show, using a century of U.S.\ county data, that
the frequency of severe natural disasters has increased significantly
since the 1990s, a trend visible in the NOAA record itself.  These
events are not only physical shocks; they are also information events.
When a hurricane makes landfall or a drought devastates a region's
agriculture, the resulting newspaper coverage is among the primary
channels through which millions of Americans form beliefs about what
happened, why, and what it portends.

Whether that coverage attributes an extreme event to anthropogenic
climate change or frames it as natural variability is a consequential
editorial choice.  The economics literature has established that
disaster media coverage materially shapes downstream behavior: using
flood events as a natural experiment, \citet{gallagher2014} finds that
insurance take-up spikes by approximately 9\% in directly affected
areas following a flood, with non-flooded communities in the same
television media market experiencing a take-up increase roughly
one-third as large---evidence that the information channel, not just
direct experience, drives risk updating.  On the policy side,
\citet{magontier2020} finds that a one-standard-deviation increase in
storm media coverage is associated with 54\% more FEMA hazard
mitigation projects in the affected area, consistent with coverage
functioning as a risk signal to governments as well as households.  If
the \emph{content} of coverage matters---if framing a disaster as a
symptom of climate change rather than as a one-off natural event shapes
how readers and policymakers update their beliefs about future risk---
then the question of what U.S.\ newspapers actually say about
causation is one with real stakes for insurance markets, preparedness
investment, and public understanding of climate risk.

Yet we know surprisingly little about the systematic pattern of causal
attribution framing across U.S.\ disaster coverage.  Prior work has
examined how coverage of climate change broadly has polarized over
time---\citet{chinn2020} show that partisan divergence in climate
discourse language was flat through most of the 2000s before a discrete
structural break around 2010--2011---and community-level studies have
documented that attribution discussion is minimal even after high-impact
events \citep{boudet2020}.  But no prior study has produced a
systematic, quantitative measure of how causal attribution framing
varies across the full universe of major U.S.\ disaster events, across
outlet types, and over time.  The barriers are practical: hand-coding
attribution framing for hundreds of events and thousands of articles is
prohibitively costly.

This paper removes that barrier.  We construct an Attribution Index
covering 1,252 newspaper articles retrieved for all 313 NOAA
billion-dollar disaster events from 2000 to 2024, using a validated
large language model (LLM) coding pipeline that follows the three-step
text-as-data framework of \citet{gentzkow2019}, including the manual
audit they recommend to cross-check fitted classifications against
human codes.  Each article is first
screened for relevance and then assigned to one of five attribution
categories: explicit anthropogenic framing ($+1$), hedged or contextual
climate mention ($+0.5$), no causal attribution ($0$), explicit
non-climate attribution ($-1$), or mixed/contested framing ($0$,
separately flagged).  The pipeline is validated against an independent
human post-hoc audit of 100 articles---blind to the LLM's codes---
which achieves $\kappa = 0.96$ for relevance screening and weighted
$\kappa_w = 0.96$--$0.98$ for attribution coding, comfortably
exceeding the reliability targets set in advance.  Article-level
scores are averaged to a disaster-level Attribution Index, enabling
systematic regression analysis of the determinants of attribution
framing at both the disaster and article level.

We document three main findings.  First, the dominant mode of disaster
coverage is silence on causation: 74.8\% of relevant articles contain
no causal attribution at all, and the mean disaster-level Attribution
Index is only 0.028---near-neutral, and far from the confident
anthropogenic framing ($+1$) that attribution science might lead one
to expect.  Second, the level of attribution framing has risen
steadily over the sample period at approximately 0.013 index points
per year ($p < 0.001$), shifting from a pre-2010 mean of $-0.105$
(net non-climate framing) to a post-2010 mean of $+0.057$ (modest net
anthropogenic framing).  This timing aligns with the structural break
in partisan climate discourse identified by \citet{chinn2020}.  Third,
outlet political lean is a robust predictor of attribution framing:
articles from more right-leaning outlets (by Ad Fontes Bias score)
receive lower Attribution Index values ($\hat\beta = -0.011$ per
Bias-score point, $p = 0.013$, clustered standard errors), a result
that survives disaster fixed effects, alternative lean measures, and an
ordered-logit functional-form check, and is concentrated in national
and wire-service coverage.  By contrast, none of the disaster's own
physical characteristics---type, CPI-adjusted cost, death toll,
duration, or geography---predict attribution framing.

The paper proceeds as follows.  Section~2 reviews the related
literature.  Section~\ref{sec:data} describes the NOAA disaster
dataset, the newspaper collection pipeline, and the construction of the
Attribution Index.  Section~\ref{sec:methods} details the LLM coding
procedure, validation design, and regression specifications.
Section~\ref{sec:results} presents the main results, and
Section~\ref{sec:results_robustness} reports the robustness checks.
Sections~7 and~8 discuss implications and conclude.


% ===================================================================
\section{Related Literature}
% ===================================================================

This paper connects four strands of research: (i)~the growing literature
on climate communication and media framing of extreme weather; (ii)~the
economics of natural disasters, salience, and risk perception;
(iii)~text-as-data and natural language processing methods in economics;
and (iv)~the political economy of local media markets.

\subsection*{Climate Communication and Media Framing}

A central concern in climate communication research is whether and how
media coverage links extreme weather events to anthropogenic climate
change.  \citet{boudet2020} study community-level responses to 15
high-impact weather events across the United States, analyzing over
4,610 newspaper articles and 164 community interviews.  They find that
local media discussion of climate change is minimal even in the
aftermath of events with a strong scientific attribution signal ---
fewer than five attribution-related articles per community on average.
Crucially, when attribution coverage does appear, both partisanship
and event-type confidence shape the discussion.  The paper distinguishes
three tiers of scientific confidence in event attribution: zero
confidence for extratropical cyclones, tornadoes, and mudslides;
intermediate confidence for droughts, wildfires, flooding, and tropical
cyclones; and high confidence for extreme heat and cold events
specifically.  Even high-confidence events do not reliably generate
sustained attribution discourse at the local level.  Our paper extends
this agenda from the small-$N$ qualitative approach of
\citet{boudet2020} to a large-$N$ quantitative analysis, trading
community-interview depth for systematic coverage of all 313 NOAA
billion-dollar disaster events over a 25-year panel.

\citet{boykoff2004} provide the foundational empirical account of how
professional norms distort climate coverage in the U.S.\ prestige press.
Analyzing 636 articles drawn from a sample of 3,543 in the \emph{New
York Times}, \emph{Washington Post}, \emph{Los Angeles Times}, and
\emph{Wall Street Journal} over 1988--2002, they find that the
journalistic norm of ``balance'' caused roughly equal column space to be
given to the scientific consensus and a small minority of skeptics:
52.65\% of articles balanced anthropogenic and natural-variability
claims, while only 35.29\% emphasized the anthropogenic view --- despite
the scientific community's growing consensus that human activities are a
primary driver of warming.  They term this ``balance as bias,'' an
informational distortion that systematically amplifies the minority
skeptical view.  The pre-2010 mean Attribution Index in our data
($-0.105$) is consistent with this balance norm persisting into the early
2000s before eroding over the subsequent decade.

\citet{boykoff2007} documents the subsequent erosion of this norm.
Extending the analysis to 2003--2006 and adding three UK newspapers
(the \emph{Guardian}/\emph{Observer}, \emph{Independent}, and \emph{The
Times}) alongside the same four U.S. prestige papers, Boykoff finds that
the share of U.S. coverage providing ``balanced accounts'' fell
from 36.6\% in 2003 (statistically significant divergence from scientific
consensus, $z = 7.73$, $p<0.001$) to just 3.3\% in 2006 (no longer
significant), while coverage attributing significant human contributions
rose from 61.0\% in 2003 to 96.7\% by 2006.  UK coverage showed no
significant divergence from scientific consensus in any year.  The study
concludes that the journalistic balance norm in U.S. prestige-press
coverage of climate change had been effectively abandoned by 2005--2006.
This erosion provides a bridge between the prior literature and the
period we study: the pre-2010 decline in biased balance sets the
baseline from which our Attribution Index begins its upward trend.

\citet{feldman2017} examine the same set of prestige newspapers---the
\emph{New York Times}, \emph{Wall Street Journal}, \emph{Washington
Post}, and \emph{USA Today}---over 2006--2011, applying a
threat-and-efficacy framework to 642 articles.  That study finds
systematic differences in framing across outlets: the \emph{Wall Street
Journal} is significantly less likely than the other three papers to
discuss climate impacts (21.6\% of articles vs.\ 40--58\% for the other
three), least likely to represent climate change as a present-day or
U.S.-specific threat, and most likely to use conflict framing (53\% of
articles) and negative economic framing (35.1\%) when discussing climate
actions.  Positive efficacy information is similarly rare across all four
papers.  This prior evidence of systematic outlet-level divergence in
prestige-press climate framing directly motivates our outlet-lean
analysis in Section~\ref{sec:outlet_results}.

\citet{feldman2012} document analogous outlet-level divergence in cable
news.  A content analysis of 269 transcripts mentioning global warming
from Fox News ($n=182$), CNN ($n=66$), and MSNBC ($n=21$) over 2007--2008
finds that approximately 60\% of Fox coverage is dismissive of climate
change versus under 20\% accepting, while CNN and MSNBC both exceed
70\% accepting.  A linked survey of a nationally representative sample
($N = 2{,}057$) finds that Fox News viewership is associated with lower
global warming acceptance ($\beta = -0.10$, $p<.001$), while CNN/MSNBC
viewership predicts higher acceptance ($\beta = +0.08$, $p<.001$) ---
effect sizes comparable to those of political partisanship.  Together
with \citet{feldman2017}, these results establish that outlet-level
differences in climate framing pervade both newspaper and broadcast
news, and are detectable in audience beliefs as well as content.

\citet{chinn2020} analyze U.S.\ climate news coverage over 1985--2017
using Wordfish, an unsupervised text-scaling method, to measure the
ideological positioning of climate discourse over time.  They find that
politicization of climate news rose gradually throughout the period,
but that partisan polarization in coverage language was relatively flat
until a discrete structural break in 2010--2011 --- precisely the
pre-/post-2010 break visible in our Figure~\ref{fig:trend_over_time}
and confirmed by our year-trend regression.  Their key outcome is
\emph{discourse polarization} --- the divergence between outlets in the
partisan language used to discuss climate --- while ours is \emph{causal
attribution framing} --- whether coverage attributes a specific event to
climate change or to natural variability.  The two complement each
other: polarization in language is a necessary but not sufficient
condition for divergence in causal attribution.  \citet{chinn2020} also
document that political actors outnumber scientists as quoted sources in
climate coverage from 2006 onward, in line with the supply-side
political-economy mechanism we revisit in the outlet-lean analysis
(Section~\ref{sec:outlet_results}).

\citet{carmichael2017} examine the drivers of public concern about
climate change using quarterly time-series data over 2001--2014.
Contrary to a simple salience hypothesis, they find that extreme weather
events do not significantly increase aggregate climate concern.  Instead,
media attention and political elite cues are the dominant drivers of
public opinion dynamics.  They document echo chamber and boomerang
effects in which partisan media exposure can \emph{decrease} concern
among ideologically unsympathetic audiences.  Our outlet-lean finding is
consistent with this mechanism at the supply side: if right-leaning
outlets systematically produce less anthropogenic attribution framing,
ideologically sorted audiences are exposed to less such framing,
reinforcing the partisan gap in concern identified by
\citet{carmichael2017}.

\citet{mccright2011} provide the audience-level foundation for this
partisan divide.  Analyzing ten nationally representative Gallup polls
(2001--2010; pooled $N$ up to 9{,}113 respondents), they document
growing ideological and partisan polarization in Americans' climate
change beliefs.  Liberals and Democrats are significantly more likely
than conservatives and Republicans to accept the scientific consensus
(e.g., 68.8\% of liberals vs.\ 42.6\% of conservatives believe global
warming effects have already begun, in pooled data), and this gap
widened substantially over the decade: the ideological spread grew from
18 percentage points in 2001 (67.1\% of liberals vs.\ 49.4\% of
conservatives) to 44 percentage points in 2010 (74.8\% vs.\ 30.2\%);
the analogous partisan gap grew from 11 to 41 percentage points.
Political orientation also moderates the effects of education and
self-reported understanding: greater knowledge and education increase
climate concern among liberals and Democrats, but the same effects are
weaker or even negative among conservatives and Republicans --- a
pattern that fits biased information processing, in which new
climate information is filtered through pre-existing partisan
commitments rather than updating beliefs uniformly.  All eight
political orientation $\times$ year interaction terms are statistically
significant ($p < .05$ to $p < .001$), confirming that this polarization
was not static but actively growing over the period our Attribution
Index covers.  This audience-level polarization is directly relevant to
the demand-side interpretation of our outlet-lean finding in
Section~\ref{sec:outlet_results}.

\citet{hai2022} use a preregistered two-wave survey experiment
(approximately 9,174 respondents total) to examine how politicians'
public attribution of extreme weather to climate change affects public
perception.  They find a backfire effect among Republicans, who rate
attribution-making politicians as \emph{less} capable.  At the same
time, they document a modest increase over time in the frequency with
which politicians make explicit attribution claims.  These results imply
a political-economy constraint on attribution language: political actors
face differential electoral incentives to link extreme events to climate
change, and those incentives vary by party and constituency.  This
mechanism echoes the partisan outlet-lean gradient we
document at the media level.

\subsection*{Economics of Natural Disasters, Salience, and Policy}

Three papers from the economics literature speak directly to how disaster
coverage translates into behavioral and policy responses.
\citet{gallagher2014} provides the most direct evidence of a
media-salience channel, showing that flood insurance take-up spikes by
approximately 9\% following a flood event in the affected area, then
decays over a nine-year window.  Critically, non-flooded communities in
the \emph{same} television media market experience a take-up increase
one-third as large as directly affected communities --- clean evidence of
a media-driven information transmission effect that operates through
coverage rather than direct experience.  Our paper measures the
climate-attribution content of exactly the kind of coverage that
\citet{gallagher2014} shows is behaviorally consequential.

\citet{boustan2020} document the macroeconomic consequences of natural
disasters using a century of U.S.\ county data.  They find that severe
disasters increase out-migration by 1.5 percentage points and depress
housing prices by 2.5--5\%; they further show that disaster frequency
has accelerated since the 1990s.  This acceleration provides the
economic backdrop for our sample: the disasters we study are a growing
risk, making the question of how media covers their causes increasingly
consequential.

\citet{magontier2020} provides the most direct empirical link between
disaster media coverage and government disaster-preparedness behavior.
Using variation in storm media coverage across U.S.\ markets, he finds
that a one-standard-deviation increase in storm media coverage is
associated with 54\% more FEMA hazard mitigation projects in the
affected area.  He proposes a risk-signaling mechanism in which media
coverage serves as a credible signal of disaster risk to voters and
policymakers, driving preparedness investment.  If the \emph{content}
of coverage matters --- whether a disaster is framed as a consequence of
climate change or as a natural event --- then the attribution framing
we measure may modulate the character of this policy response; this is
a question our paper motivates but does not directly
answer.\footnote{\citet{magontier2020} is currently a working paper;
its empirical findings should be interpreted with appropriate caution
pending peer review.}

\citet{egan2012} provide a complementary micro-level result at the
individual level.  Exploiting within-person variation in local
temperature departures from seasonal norms, they find that a
3.1\textdegree F increase in local temperature above its seasonal norm
raises the probability that a respondent says there is ``solid
evidence'' of warming by approximately 1 percentage point, with the
effect decaying within 12 days.  The effect is largest among respondents
who are less educated and who have weak attachments to political parties.
This evidence that direct environmental experience shapes climate beliefs
through a short-term salience mechanism complements our evidence that
media attribution framing shapes the broader information environment
in which those experiences are interpreted.

\subsection*{Text as Data and NLP Methods in Economics}

The methodological approach of this paper --- using a large language
model to classify natural language text at scale --- sits within an
expanding economics literature on text-as-data methods.
\citet{gentzkow2019} provide a comprehensive survey of this literature,
organizing it around a three-step framework: (i)~represent text as a
numerical array (via bag-of-words, embeddings, or model outputs);
(ii)~map those representations to predicted values or classifications;
and (iii)~use the predicted values in subsequent economic analysis.
Our pipeline maps directly onto this framework, with steps (i)--(ii)
implemented by a large language model (Claude, Anthropic) following a
written codebook, and step (iii) implemented by the regression analysis
in Section~\ref{sec:results}.  \citet{gentzkow2019} also specifically
recommend manual audits --- cross-checking a subset of fitted values
against human-produced codes --- as a best practice for validating
text-as-data pipelines; our independent 100-article post-hoc audit
(Section~\ref{sec:results_robustness}) follows this recommendation
directly.

\citet{gentzkow2010} represent the paradigmatic application of
text-as-data methods to measure media bias, using the frequency of
partisan phrases in congressional speech to construct a slant index for
over 400 U.S.\ daily newspapers.  They find that approximately 20\% of
the variation in outlet slant is explained by reader demographic demand,
while owner identity explains far less.  This result disciplines our
interpretation of the outlet-lean gradient we document: partisan framing
differences across outlets appear to reflect audience demand as much as
proprietorial preferences.  Their approach uses the text of newspapers
as the object of analysis, exactly as we do, though they measure
partisan language use while we measure causal attribution framing.

\subsection*{Political Economy of Local Media Markets}

Our finding that the political-lean effect is concentrated in
national-outlet coverage raises questions about the role of local media
market structure more broadly.  \citet{ellger2024} study the causal
effect of local newspaper closures on political polarization using
German newspaper exits over 1980--2009, exploiting exit events as
quasi-experimental variation in local newspaper supply.  Using a
difference-in-differences design, they find that local newspaper exits
increase political polarization, and that the mechanism operates through
substitution: when a local paper exits, consumers shift to national
tabloid coverage that is more nationally partisan and less locally
focused.  This result offers a candidate explanation for the pattern
we document --- local-outlet attribution framing shows no significant
lean gradient (Section~\ref{sec:outlet_results}), while national-outlet
framing shows a clear partisan pattern --- if local papers play a
distinct informational role that is not replicable by national
substitutes, their ongoing decline could gradually shift the aggregate
information environment toward the nationally partisan framing that
dominates in our national-outlet subsample.


% ===================================================================
\section{Data}
\label{sec:data}
% ===================================================================

\subsection{NOAA Billion-Dollar Disasters Dataset}

The sample is drawn from NOAA's Weather and Climate Billion-Dollar
Disasters dataset, restricted to U.S.\ events occurring between 2000
and 2024 (313 events). Each event is characterized by a disaster type
(Severe Storm, Tropical Cyclone, Flooding, Drought, Wildfire, Winter
Storm, or Freeze), a begin and end date, the set of affected states,
CPI-adjusted and unadjusted cost (in millions of 2024 dollars), and a
fatality count. Severe storms are the most common event type (179 of
313 events), followed by tropical cyclones (48), flooding (30),
drought (21), wildfire (19), winter storm (13), and freeze (3). Total
CPI-adjusted damages across the sample exceed \$2.36 trillion, with a
median event cost of \$2.4 billion (mean \$7.5 billion, driven by a
right tail of extreme events up to \$201.3 billion); total reported
fatalities across all events are 10,849.

\subsection{Newspaper Coverage Data}

\paragraph{Newspaper universe.} The candidate universe of outlets is
drawn from the Access World News (AWN) database, which lists 3,647
U.S.\ newspapers and wire services with associated city, state, and
language metadata. Of these, 970 had no listed website and were
excluded. The remaining outlets were passed through an automated
verification pipeline that attempted to resolve and HTTP-check each
listed URL (classifying each as verified, verified via HTTP-only,
or unverified with a status code/reason); 899 outlets ultimately
produced a usable, live domain.

\paragraph{Programmable Search Engine (PSE) construction.} The 899
verified domains were partitioned into 18 regional Google
Programmable Search Engines (PSEs), grouped by the states each outlet
primarily serves, plus a 19th ``National'' PSE covering major national
and wire outlets (e.g., the \emph{New York Times}, \emph{Washington
Post}, CNN, AP, NPR, USA Today) that are relevant regardless of event
geography. Each of the 313 disaster events is mapped to the subset of
regional PSEs corresponding to its affected states (plus the National
PSE), so that searches are restricted to outlets plausibly covering
that event's geography.

\paragraph{Collection pipeline.} For each event, queries combining the
disaster name/type and date window are issued against the relevant
PSE subset via the Google Custom Search JSON API. Candidate results
are filtered on (i) a date window around the event's begin/end dates
(with retrospective coverage allowed up to the start of the next event
of the same type, where applicable), (ii) geographic relevance --
results must reference a place name consistent with the event's
affected states or broader region, and (iii) disaster-type/keyword
relevance. Up to 30 candidate articles are retained per event. The
full collection run returned 1,694 candidate articles across the 313
events: 237 events (75.7\%) returned at least one candidate article
(mean 5.4 articles per event among events with any coverage; 76
events returned none). Of the 1,694 candidate articles, 911 (53.8\%)
are from outlets on the National PSE list and 783 (46.2\%) are from
regional/local outlets. These candidate articles are then passed
through the relevance-screening and attribution-coding procedure
described in Section~\ref{sec:methods}.

\subsection{Constructing the Attribution Index}
\label{sec:attribution_index_data}

Each of the 1,694 candidate articles is first screened for
\emph{relevance} -- whether it substantively discusses the specific
NOAA-listed event it was retrieved for (correct disaster, geography,
and time window), versus being off-topic, ambiguous (\texttt{Borderline}),
or unusable (\texttt{Trash}). Only articles coded \texttt{Relevant}
proceed to attribution coding.

Each relevant article is then assigned to one of five categories
describing how it frames the causal origin of the event, with an
associated numeric score $a_{i,d} \in \{-1, 0, 0.5, 1\}$:

\begin{table}[ht]
\centering
\caption{Attribution Coding Categories}
\label{tab:attribution_categories}
\begin{threeparttable}
\begin{tabular}{cp{10.5cm}c}
\toprule
Code & Description & Score \\
\midrule
A & Explicit anthropogenic climate-change attribution & $+1.0$ \\
B & Hedged / contextual climate mention & $+0.5$ \\
C & No causal attribution (pure event reporting) & $0.0$ \\
D & Explicit non-climate attribution (natural variability or other human causes) & $-1.0$ \\
E & Mixed / contested framing & $0.0$ (flagged) \\
\bottomrule
\end{tabular}
\begin{tablenotes}
\small
\item Notes: Full category definitions, decision rules, and example
language are given in the project codebook
(\texttt{paper/attribution\_codebook.md}). Category E receives the
same numeric score as C but is separately flagged as
\emph{contested} when computing disaster-level dispersion.
\end{tablenotes}
\end{threeparttable}
\end{table}

Article-level scores are averaged to a disaster-level
\emph{Attribution Index}; the aggregation formula, calibration
procedure, and treatment of events with no relevant coverage are
described in Section~\ref{sec:methods}.

\subsection{Summary Statistics}
\label{sec:summary_stats}

Table~\ref{tab:summary_stats} summarizes the collection-stage data
described above, together with the results of the relevance-screening
and attribution-coding pass described in Section~\ref{sec:methods}.
Of the 1,694 candidate articles, full text was successfully retrieved
for 1,591 (93.9\%); of those, 1,252 (78.7\%) were coded
\texttt{Relevant} to the specific event they were retrieved for, with
the remainder split between \texttt{Borderline} (210, 13.2\%) and
\texttt{Trash} (129, 8.1\%). Among Relevant articles, the large
majority (937, 74.8\%) contain no causal attribution at all (category
C, pure event reporting); explicit anthropogenic attribution (A) and
hedged/contextual climate mentions (B) together account for 252
articles (20.1\%), explicit non-climate attribution (D) for 53
(4.2\%), and mixed/contested framing (E) for only 10 (0.8\%) --
contested coverage is rare both at the article level and at the
disaster level (only 6 of the 209 covered disasters, 2.9\%, have any
Contested article at all). At the disaster level, 209 of the 313
events (66.8\%) have at least one Relevant article (mean 6.0, median
4 Relevant articles per covered event); the remaining 104 (33.2\%)
have no Relevant coverage at all and are analyzed as a separate
no-attribution group rather than dropped (Section~\ref{sec:methods}).
The mean disaster-level Attribution Index among covered events is
0.028 (median 0.000, SD 0.241, range $[-1, 1]$) -- i.e.\ coverage is,
on average, only very mildly tilted toward anthropogenic framing and
is dominated in practice by events with little to no causal framing
either way, consistent with the category-C share above.

\begin{table}[ht]
\centering
\caption{Summary Statistics: Sample, Coverage, and Coding Outcomes}
\label{tab:summary_stats}
\begin{threeparttable}
\begin{tabular}{lc}
\toprule
 & Value \\
\midrule
NOAA billion-dollar disaster events (2000--2024) & 313 \\
Candidate newspaper articles collected & 1,694 \\
Events with $\geq 1$ candidate article & 237 (75.7\%) \\
Events with no candidate articles & 76 (24.3\%) \\
Mean articles per covered event & 5.4 \\
Maximum articles for a single event & 30 \\
Articles from National PSE outlets & 911 (53.8\%) \\
Articles from regional/local PSE outlets & 783 (46.2\%) \\
\midrule
\multicolumn{2}{l}{\emph{Full-text retrieval and coding outcomes:}} \\
Candidate articles with full text retrieved (``found'') & 1,591 (93.9\%) \\
Articles coded \texttt{Relevant} (of found) & 1,252 (78.7\%) \\
Articles coded \texttt{Borderline} (of found) & 210 (13.2\%) \\
Articles coded \texttt{Trash} (of found) & 129 (8.1\%) \\
\quad Category A -- explicit anthropogenic (of Relevant) & 107 (8.5\%) \\
\quad Category B -- hedged/contextual climate mention & 145 (11.6\%) \\
\quad Category C -- no causal attribution & 937 (74.8\%) \\
\quad Category D -- explicit non-climate attribution & 53 (4.2\%) \\
\quad Category E -- mixed/contested & 10 (0.8\%) \\
\midrule
\multicolumn{2}{l}{\emph{Disaster-level Attribution Index:}} \\
Disasters with $\geq 1$ Relevant article ($N_d > 0$) & 209 (66.8\%) \\
Disasters with no Relevant coverage ($N_d = 0$) & 104 (33.2\%) \\
Mean / median $N_d$ among covered disasters & 6.0 / 4 \\
Mean $\text{AttributionIndex}_d$ (covered disasters) & 0.028 \\
Median / SD $\text{AttributionIndex}_d$ & 0.000 / 0.241 \\
Disasters with any Contested coverage & 6 (2.9\% of covered) \\
\bottomrule
\end{tabular}
\begin{tablenotes}
\small
\item Notes: Coverage statistics in the top panel describe the raw
candidate-article sample prior to relevance screening. The bottom two
panels reflect the full-sample LLM coding pass described in
Section~\ref{sec:methods}; reliability statistics from the
independent post-hoc audit sample are reported separately once
available (Section~\ref{sec:methods}).
\end{tablenotes}
\end{threeparttable}
\end{table}


% ===================================================================
\section{Empirical Approach}
\label{sec:methods}
% ===================================================================

\subsection{Article-Level Classification}

Each candidate article is coded in two sequential steps, following a
written codebook (reproduced in full in the project repository as
\texttt{attribution\_codebook.md}) that is also used as the basis for
the LLM coding prompt.

\paragraph{Step 1: Relevance screening.} Every candidate article is
classified as \texttt{Relevant} (substantively discusses the specific
NOAA-listed event -- correct disaster, geography, and time window),
\texttt{Borderline} (plausibly related but ambiguous, e.g.\ a brief
wire item or a retrospective mention), or \texttt{Trash} (off-topic,
wrong event, duplicate, paywall/landing page, dead link, or
photo gallery with caption-only text). Only
\texttt{Relevant} articles proceed to Step 2; \texttt{Borderline} and
\texttt{Trash} articles are retained in the data for transparency but
excluded from the index.

\paragraph{Step 2: Attribution coding.} Each \texttt{Relevant} article
is assigned to one of five categories describing how it frames the
causal origin of the event (Table~\ref{tab:attribution_categories}):
(A) explicit anthropogenic attribution ($+1.0$), (B) hedged or
contextual climate mention ($+0.5$), (C) no causal attribution
($0.0$), (D) explicit non-climate attribution to natural variability
or other human causes ($-1.0$), and (E) mixed/contested framing
($0.0$, separately flagged as \texttt{Contested}). For category D,
coders additionally record whether the non-climate attribution is to
natural variability (``just weather'') or to other human causes
(e.g.\ land use, infrastructure), since both score $-1$ on the
climate-vs-natural axis but imply different policy narratives.

\paragraph{Full article text retrieval.} All coding---both the human
calibration pass and the LLM full-sample pass---is performed on the
complete article text rather than search-engine snippets. Full text
was retrieved in two stages. First, archived copies of each article
were obtained via archive.ph, a web archiving service that captures
complete page content prior to paywall overlays. For articles
unavailable via archive.ph, full text was retrieved manually via
Factiva (Dow Jones). This two-stage retrieval procedure was adopted
following a finding during calibration that critical attribution
language can appear mid-article and is missed with snippet-only
coding. Of the 1,694 candidate articles, full text was ultimately
retrieved for 1,591 (93.9\%); the remaining 103 are excluded from LLM
coding and treated as missing in the Attribution Index, comprising 64
articles for which full text could not be recovered by either method
(e.g.\ dead links, HTTP errors, or pages with no archived snapshot),
36 articles that resolved to photo/video-only gallery pages with no
substantive article text, and 3 articles -- all from
\texttt{transcripts.cnn.com} -- whose URL resolved to a generic
chronological listing of segments rather than the specific transcript
segment retrieved for that event.

\paragraph{Calibration and validation.} Coding is performed by an
LLM following the codebook. The LLM is Claude (model
\texttt{claude-sonnet-4-6}; Anthropic, Inc.), accessed via the
Anthropic Messages API, with the codebook supplied as a system prompt
and the full article text as the user message. Validation follows a
staged procedure that applies two different reliability bars: a low
gate during calibration to catch codebook problems cheaply, and a
higher target for the reliability statistics reported in the paper.
First, the researcher hand-codes a stratified random calibration
sample of 50 articles (stratified by disaster type and by
pre-/post-2010 period) for both Relevance and Attribution category.
Second, the LLM codes the same 50 articles from the codebook alone.
Third, human and LLM codes are compared using unweighted Cohen's
$\kappa$ for the categorical Relevance (3-category) and Attribution
(5-category) codes; if agreement is weak ($\kappa < 0.6$,
``moderate'' or worse on the \citet{landiskoch1977} scale), the
codebook wording and examples are revised and the comparison is
repeated before scaling up. This 0.6 threshold is intentionally
low---its purpose is to catch broken category definitions early, not
to certify reliability.

Agreement on the initial calibration pass exceeded the 0.6 threshold for both codes. Unweighted
$\kappa$ for Relevance was 0.81 (``almost perfect agreement'' on the
\citealt{landiskoch1977} scale, exceeding the $\kappa \geq 0.7$
reporting target). For Attribution, both unweighted and linearly
weighted $\kappa$ were 1.00 (perfect agreement), with a mean absolute
difference of 0.00 between human- and LLM-assigned numeric scores.
The four residual Relevance disagreements were all boundary cases
(\texttt{Trash} vs.\ \texttt{Borderline}) where the correct
classification is genuinely ambiguous; the relevant decision rules
have been noted in the codebook for the full-sample pass.

Once this gate was cleared, the LLM coded all remaining articles in
the full sample (1,591 articles with retrievable full text; the 50
calibration-sample articles were not re-coded by the full pass and
instead carry forward their calibration-stage LLM codes, for
consistency of method across the full \texttt{Master} sheet).

Finally, an independent post-hoc audit sample is drawn from the fully
coded data, excluded from the calibration sample, for a final human
spot-check. Because attribution categories D and E are rare in the
full corpus (4.2\% and 0.8\% of Relevant articles respectively), a
simple random draw of the target 100 articles would likely contain
too few instances of either to support a meaningful reliability
statistic for them. The audit sample therefore deliberately
oversamples these categories -- including all 10 category-E articles
and a random 20 of the 53 category-D articles -- with the remaining
70 slots filled by a proportionally stratified random draw (by
disaster type and pre-/post-2010 period) from the rest of the coded
corpus, for a total of 100 articles. Because of this deliberate
oversampling, the audit sample is not a simple random sample of the
corpus; reliability statistics computed from it should be read as
conditional on this design (informative about each category's
reliability individually) rather than as a corpus-representative
error rate, and this is noted explicitly when those statistics are
reported. The audit sample is hand-coded blind to the LLM's codes and
then compared to them. For Relevance we report unweighted Cohen's
$\kappa$ (target $\kappa \geq 0.7$); for Attribution, we report both
unweighted $\kappa$ and a weighted $\kappa$ (linear weights) that
reflects the ordinal structure of the underlying scores
$\{-1, 0, 0.5, 1\}$, since unweighted $\kappa$ treats an A-vs-D
disagreement the same as a B-vs-C disagreement despite the very
different implied score gap (target $\kappa_w \geq 0.7$, ideally
$\geq 0.8$); and we additionally report the mean absolute difference
between human- and LLM-assigned numeric scores as a directly
interpretable companion statistic. Should the audit-sample statistics
fall short of these targets, this is disclosed and discussed as a
limitation, since the boundaries between categories B/C and C/E are
genuinely fuzzy even for human coders.

Results from the completed audit sample are as follows. For Relevance (all 100 audited
articles), raw agreement was 99\% and unweighted Cohen's $\kappa = 0.96$,
comfortably exceeding the $\kappa \geq 0.7$ target; the single
disagreement was a Trash/Borderline boundary case. For Attribution,
restricted to the 86 articles both the human coder and the LLM
independently classified as Relevant (the subsample where both raters
have a meaningful 5-category code), raw agreement was 95.3\%, unweighted
$\kappa = 0.93$, linearly weighted $\kappa_w = 0.96$, and quadratically
weighted $\kappa_w = 0.98$ -- all well above the $\kappa \geq 0.7$ target
and the $\kappa_w \geq 0.8$ ``ideal'' threshold. The mean absolute
difference between human- and LLM-assigned numeric scores was $0.023$ on
the $[-1, 1]$ scale. All four Attribution disagreements fell on adjacent
categories on the causal-framing spectrum (two A/B, one B/C, one B/E) --
exactly the kind of boundary fuzziness anticipated above -- with zero
disagreements spanning a larger gap (e.g.\ no A-vs-D confusions). Full
counts and the complete confusion matrices are reproducible from
\texttt{code/audit\_reliability.py}.

\subsection{Disaster-Level Attribution Index}

For each disaster $d$ with $N_d$ articles coded \texttt{Relevant},
article-level scores are averaged to a disaster-level
\emph{Attribution Index}:

\begin{equation}
\text{AttributionIndex}_{d} = \frac{1}{N_d}\sum_{i=1}^{N_d} a_{i,d},
\qquad a_{i,d} \in \{-1,\ 0,\ 0.5,\ 1\}
\label{eq:attribution_index}
\end{equation}

where $a_{i,d}$ is the numeric score for article $i$ covering disaster
$d$ from Table~\ref{tab:attribution_categories}, with $-1$ indicating
explicit non-climate (natural/random or other human-cause) attribution
and $+1$ indicating explicit anthropogenic climate-change attribution.
Category E (mixed/contested) scores $0$, identically to category C (no
attribution), but is additionally flagged so that a disaster-level
\emph{Contested share} -- the fraction of $N_d$ articles coded E -- can
be reported separately as a measure of framing disagreement, distinct
from the simple absence of any causal claim.

For the 104 events with $N_d = 0$ (no relevant articles, including
events that returned zero search results entirely),
$\text{AttributionIndex}_d$ is left undefined (missing) rather than
imputed, and these events are analyzed as a separate ``no media
attribution'' group rather than dropped from the sample. All of the
above quantities -- $N_d$, $\text{AttributionIndex}_d$, and the
Contested share -- are computed via formulas in the \texttt{All
Disasters} sheet of the project workbook (\texttt{attribution.xlsx}),
which references the article-level codes in the \texttt{Master}
sheet.

\subsection{Specification}
\label{sec:specification}

We relate the Attribution Index to disaster and outlet characteristics at
two levels, deliberately. Disaster type, severity, geography, and time are
properties of the \emph{event} and do not vary across the articles covering
it, so they are tested in a disaster-level cross-section. National/local
status and outlet political lean, by contrast, vary \emph{within} a
disaster -- the same wildfire is often covered by both a wire service and a
local paper -- so collapsing them to one disaster-level number would
discard most of the identifying variation; these are instead tested in an
article-level specification with standard errors clustered by disaster.

\paragraph{Disaster-level specification.} For each disaster $d$ with
$N_d \geq 1$ (i.e.\ excluding the 104 events with no Relevant coverage):

\begin{multline}
\text{AttributionIndex}_{d} = \beta_0 + \boldsymbol{\delta}'\,\text{Type}_d
  + \beta_1 \ln(\text{Cost}_d) + \beta_2 \ln(\text{Deaths}_d + 1)
  + \beta_3 \ln(\text{Duration}_d) \\
  + \beta_4 \, \text{CoastalShare}_d
  + \beta_5 \, \text{RedShare}_d + \beta_6 \, \text{Year}_d
  + \beta_7 \, N_d + \varepsilon_d
\label{eq:main_spec}
\end{multline}

where $\text{Type}_d$ is a vector of disaster-type dummies (reference
category: Severe Storm, the modal type); cost and deaths are
log-transformed given their strong right skew ($\ln(\text{Deaths}+1)$ to
accommodate zero-fatality events); $\ln(\text{Duration}_d)$ is the log of
the number of calendar days between the disaster's NOAA-listed begin and
end dates; $\text{CoastalShare}_d$ and $\text{RedShare}_d$ are,
respectively, the share of $d$'s affected states that border an ocean or
the Gulf of Mexico and the share that voted Republican in the 2020
presidential election -- shares rather than a single ``primary state''
because the \texttt{States} field lists affected states alphabetically, not
by impact, so no reliable primary state can be recovered for multi-state
events; $\text{Year}_d$ is the calendar year of the disaster's begin date;
and $N_d$ (the number of Relevant articles) is included as a control for
coverage volume. Standard errors are heteroskedasticity-robust (HC1).

\paragraph{Article-level specification.} For each Relevant article $i$
covering disaster $d$:

\begin{equation}
a_{i,d} = \gamma_0 + \gamma_1 \, \text{National}_{i} + \gamma_2 \,
  \text{OutletLean}_{i} + \boldsymbol{\delta}'\,\text{Type}_d
  + \gamma_3 \, \text{Year}_d + u_{i,d}
\label{eq:article_spec}
\end{equation}

where $\text{National}_i$ is the existing \texttt{National Newspaper}
indicator and $\text{OutletLean}_i$ is the outlet's Ad Fontes Bias score
(continuous, $-42$ to $+42$, more positive = more right-leaning), joined by
domain; standard errors are clustered by disaster $d$ to account for
non-independence among articles covering the same event.
$\text{OutletLean}_i$ is only defined for the 33 of 135 distinct domains in
the sample with an individually-rated page on Ad Fontes/AllSides (covering
71\% of Relevant articles, since coverage is concentrated in large
outlets), so specifications including it are estimated on that subsample.


% ===================================================================
\section{Results}
\label{sec:results}
% ===================================================================

All results below are estimated on the 209 disasters with $N_d \geq 1$ (or
the relevant article-level subsample). The independent post-hoc audit
described in Section~\ref{sec:methods} clears the reliability targets
($\kappa = 0.96$ for Relevance, weighted $\kappa_w = 0.96$--$0.98$ for
Attribution); see Section~\ref{sec:results_robustness} for the full
statistics.

\subsection{Descriptive Patterns}

Figure~\ref{fig:trend_over_time} plots the disaster-level Attribution Index
against year, together with a fitted linear trend. The index is centered
near zero (or negative) for most of the sample's first decade and rises
steadily afterward: mean $\text{AttributionIndex}_d$ is $-0.105$
(SD $0.250$) for disasters beginning before 2010, versus $+0.057$
(SD $0.230$) from 2010 onward, a difference that is statistically
significant (Welch $t = -3.63$, $p < 0.001$). A bivariate regression of the
index on year gives a slope of $0.0129$ per year ($p < 0.001$,
robust SE), i.e.\ coverage has moved from net-skeptical/natural framing
toward modestly net-anthropogenic framing over the sample period, though
even at the end of the sample the average disaster's coverage remains much
closer to ``no causal claim'' (category C) than to confident anthropogenic
attribution.

\begin{figure}[ht]
\centering
\includegraphics[width=0.8\textwidth]{figures/attribution_trend.pdf}
\caption{Attribution Index Over Time. Points are annual means across
covered disasters ($\pm$1 SE); line is an OLS fit of
$\text{AttributionIndex}_d$ on year (slope $=0.0129$, $p<0.001$).}
\label{fig:trend_over_time}
\end{figure}

\FloatBarrier
Appendix Figure~\ref{fig:index_hist} shows the full cross-sectional
distribution underlying that trend. The index is exactly $0.0$ for 55.5\%
of covered disasters (consistent with category C, no causal attribution,
being the modal article-level code), with 30.6\% net-positive (mean tilt
toward anthropogenic framing) and 13.9\% net-negative; most of the non-zero
mass sits close to zero in both directions, and disasters with a strongly
net-anthropogenic ($\text{AttributionIndex}_d$ near $+1$) or net-skeptical
(near $-1$) coverage profile are rare. This reinforces the reading from
Section~\ref{sec:summary_stats}: the year-over-year movement documented
above is a shift in the modest tilt of an otherwise near-zero-centered
distribution, not a shift between two opposed extremes.

\FloatBarrier
\subsection{Variation by Geography}

Splitting disasters into ``Coastal'' (at least one affected state borders
an ocean or the Gulf of Mexico) versus ``Inland'' shows no significant
difference: mean Attribution Index is $0.023$ (SD $0.235$, $N=179$) for
Coastal disasters and $0.070$ (SD $0.301$, $N=22$) for Inland disasters
(Welch $t=-0.72$, $p=0.48$). Classifying disasters by the 2020 presidential
vote of their affected states (majority-rule across states, for disasters
spanning more than one state) likewise shows essentially no difference:
$0.023$ (SD $0.247$) for Majority-Blue disasters ($N=71$) versus $0.023$
(SD $0.239$) for Majority-Red disasters ($N=125$); the 5 Split/Tied
disasters average higher ($0.213$) but the group is too small to draw any
conclusion from. A regression of the index on the continuous
$\text{CoastalShare}_d$ and $\text{RedShare}_d$ measures (which use the
fraction of a multi-state disaster's states that are coastal/Republican,
rather than collapsing to a single category) confirms both are
statistically indistinguishable from zero (Table~\ref{tab:main_results},
columns 2--3). Eight disasters with \texttt{States = NATIONAL} or limited
to Puerto Rico (which has no 2020 presidential vote) are excluded from this
section and from any specification including geography.

\subsection{Variation by Disaster Type and Severity}

Mean Attribution Index varies little across NOAA disaster types -- from
$0.000$ (Freeze, $N=2$) to $0.074$ (Winter Storm, $N=12$) -- and a one-way
ANOVA finds no significant difference across the seven types ($F=0.32$,
$p=0.93$; full per-type means in Appendix Table~\ref{tab:type_means}, full
type-dummy regression coefficients in Appendix
Table~\ref{tab:appendix_type_full}). Severity, measured as
log-CPI-adjusted cost and log(deaths$+1$), is similarly uninformative on
its own (correlation with the index of $0.09$ and $0.07$ respectively,
neither coefficient significant at conventional levels in
Table~\ref{tab:main_results}); event duration (log days) is also not a
significant predictor. In short, none of the disaster's own physical
characteristics -- what kind of disaster it was, how costly, how deadly, or
how long it lasted -- predicts how its coverage frames causation, at least
in this disaster-level cross-section.

\subsection{Variation by Outlet Characteristics}
\label{sec:outlet_results}

Unlike disaster type, severity, and geography, outlet characteristics show
up clearly -- but only once a confound is accounted for. In a simple
bivariate comparison, articles from national outlets score modestly higher
than local outlets ($0.120$ vs.\ $0.077$ on the $[-1,1]$ scale; clustered
$t$-test $p=0.079$), but this gap is not robust to controlling for
disaster type and year: in the multivariate article-level specification
(Table~\ref{tab:outlet_results}, column 2), the National coefficient falls
to $0.031$ and is no longer statistically distinguishable from zero
($p=0.46$). National outlets in this sample skew toward later years and
toward disaster types that independently trend toward more anthropogenic
framing, and that composition effect appears to explain most of the gap in
the bivariate comparison.

Outlet political lean tells a different story. Using the Ad Fontes Bias
score (continuous, more positive = more right-leaning) for the 71\% of
Relevant articles from an individually-rated domain, a one-point increase
in Bias score is associated with a $0.011$-point \emph{decrease} in
Attribution Index ($p=0.013$, clustered SE), controlling for disaster type,
year, and national/local status (Table~\ref{tab:outlet_results}, column
2) -- i.e.\ articles from more right-leaning outlets are systematically
less likely to frame disasters as anthropogenically caused, even after
accounting for which disasters and which years they're covering. This
holds up in the bivariate specification as well ($-0.013$, $p=0.002$) and
is robust to using the AllSides score instead of Ad Fontes
($-0.033$, $p=0.003$; Section~\ref{sec:results_robustness}). Appendix
Figure~\ref{fig:outlet_lean_scatter} plots the underlying relationship: a
jittered scatter of the four discrete article-level scores against outlet
Bias score, with the OLS fit and binned means overlaid -- the binned means
track the linear fit closely, suggesting the relationship is not an
artifact of a nonlinear pattern being forced into a straight line.
(Section~\ref{sec:results_robustness} also shows this effect is
concentrated in national-outlet coverage and does not hold among
local-only articles.)

\begin{table}[ht]
\centering
\caption{Outlet Characteristics: Article-Level Results}
\label{tab:outlet_results}
\begin{threeparttable}
\begin{tabular}{lcc}
\toprule
 & (1) & (2) \\
 & \multicolumn{2}{c}{Dep.\ var.: $a_{i,d}$ (article attribution score)} \\
\midrule
National outlet      & 0.0308   & 0.0308   \\
                      & (0.0418) & (0.0418) \\
Ad Fontes Bias score  &          & $-$0.0112$^{**}$ \\
                      &          & (0.0045) \\
Year                  & 0.0169$^{***}$ & 0.0177$^{***}$ \\
                      & (0.0024) & (0.0027) \\
\midrule
Disaster type dummies & Yes & Yes \\
Clustered SE (by disaster) & Yes & Yes \\
$N$ (articles)        & 1{,}252 & 889 \\
Clusters (disasters)  & 209 & 194 \\
$R^2$                 & 0.112 & 0.124 \\
\bottomrule
\end{tabular}
\begin{tablenotes}
\small
\item Notes: $^{*}p<0.10$, $^{**}p<0.05$, $^{***}p<0.01$. Standard errors
in parentheses, clustered by disaster. Column (1) is estimated on all
1,252 Relevant articles; column (2) is restricted to the 889 articles
(71\%) from one of the 33 individually-rated domains, since Ad Fontes
lean is undefined otherwise. Full disaster-type dummy coefficients are in
Appendix Table~\ref{tab:appendix_article_full}.
\end{tablenotes}
\end{threeparttable}
\end{table}

\clearpage
\subsection{Main Regression Results}

Table~\ref{tab:main_results} reports the disaster-level specification
(equation~\ref{eq:main_spec}) in three nested columns: (1) disaster type
and severity only; (2) adding geography and duration; (3) the full
specification, adding year and coverage volume ($N_d$). Disaster type,
cost, deaths, duration, and geography are not significant in any column.
Year is strongly significant once added ($\beta=0.0128$, $p<0.001$),
consistent with Figure~\ref{fig:trend_over_time}, and coverage volume
$N_d$ is also positive and significant ($\beta=0.0089$, $p=0.036$) --
disasters that draw more press coverage skew slightly more anthropogenic
on average, though the effect is small and we would caution against a
causal reading (more coverage could itself be a consequence of a disaster
having a salient climate angle, rather than a cause of one). The full
specification explains 20.1\% of disaster-level variance in the Attribution
Index, almost entirely attributable to the year and $N_d$ terms -- the
disaster's own type, cost, death toll, and geography contribute almost
nothing on top of those.

\begin{table}[ht]
\centering
\caption{Determinants of the Attribution Index (Disaster Level)}
\label{tab:main_results}
\begin{threeparttable}
\begin{tabular}{lccc}
\toprule
 & (1) & (2) & (3) \\
 & \multicolumn{3}{c}{Dep.\ var.: $\text{AttributionIndex}_d$} \\
\midrule
$\ln(\text{Cost})$        & 0.0226   & 0.0254   & 0.0120   \\
                          & (0.0188) & (0.0203) & (0.0175) \\
$\ln(\text{Deaths}+1)$    & $-$0.0061 & $-$0.0057 & 0.0045   \\
                          & (0.0182) & (0.0191) & (0.0166) \\
$\ln(\text{Duration})$    &          & $-$0.0203 & 0.0012   \\
                          &          & (0.0214) & (0.0213) \\
CoastalShare$_d$          &          & $-$0.0191 & $-$0.0214 \\
                          &          & (0.0664) & (0.0628) \\
RedShare$_d$              &          & 0.0338   & 0.0164   \\
                          &          & (0.0553) & (0.0549) \\
Year                      &          &          & 0.0128$^{***}$ \\
                          &          &          & (0.0036) \\
$N_d$ (coverage volume)   &          &          & 0.0089$^{**}$  \\
                          &          &          & (0.0042) \\
\midrule
Disaster type dummies & Yes & Yes & Yes \\
$N$ (disasters)        & 201 & 201 & 201 \\
$R^2$                  & 0.015 & 0.019 & 0.201 \\
\bottomrule
\end{tabular}
\begin{tablenotes}
\small
\item Notes: $^{*}p<0.10$, $^{**}p<0.05$, $^{***}p<0.01$.
Heteroskedasticity-robust (HC1) standard errors in parentheses. Sample
restricted to the 201 (of 209) covered disasters with classifiable
geography (excludes 8 disasters coded \texttt{States = NATIONAL} or
Puerto Rico only). Reference disaster type is Severe Storm; full
type-dummy coefficients are in Appendix
Table~\ref{tab:appendix_type_full}.
\end{tablenotes}
\end{threeparttable}
\end{table}


% ===================================================================
\FloatBarrier
\section{Robustness}
\label{sec:results_robustness}
% ===================================================================

\paragraph{Alternative outlet-lean measure.} The outlet political-lean
result in Section~\ref{sec:outlet_results} uses Ad Fontes' numeric Bias
score because every one of the 33 individually-rated domains in the sample
has a clean numeric value there, whereas a handful of AllSides ratings
returned only a category label with no exact meter value (mapped, for this
check, to a fixed scale: Left $=-5$, Lean Left $=-2.5$, Center $=0$, Lean
Right $=+2.5$, Right $=+5$). Re-estimating the bivariate article-level
specification with the AllSides score in place of Ad Fontes gives the same
sign and a comparable magnitude ($-0.033$, $p=0.003$, clustered SE,
$N=930$), so the political-lean finding is not an artifact of which rating
service is used.

\paragraph{Disaster fixed effects.} The article-level outlet-lean
specification in Table~\ref{tab:outlet_results} controls for disaster type
and year but only \emph{clusters} standard errors by disaster, leaving
open the possibility that some other disaster-level confound drives the
result. Re-estimating with disaster fixed effects -- so that
$\text{National}_i$ and $\text{OutletLean}_i$ are identified purely from
variation \emph{within} the same disaster (e.g.\ the AP vs.\ a local paper
covering the same wildfire), implemented via within-disaster demeaning on
the 144 disasters with $\geq 2$ lean-rated articles -- leaves the lean
coefficient intact ($-0.012$, $p=0.046$) and the national coefficient
modestly larger and now marginally significant ($0.075$, $p=0.099$). This
is the most demanding test available in this sample, and the lean result
survives it.

\paragraph{Alternative attribution measures.} Three checks probe whether
the year trend (Section~\ref{sec:results}) and the outlet-lean effect
(Section~\ref{sec:outlet_results}) depend on the specific
$\{-1,0,0.5,1\}$ scoring of the Attribution Index. First, collapsing to a
binary ``any climate mention'' (A/B) vs.\ ``none'' (C/D/E) measure leaves
both findings significant (year trend $\beta=0.0100$, $p<0.001$; outlet
lean $\beta=-0.0177$, $p<0.001$). Second, recoding category-D articles
attributed to \emph{other human causes} (rather than natural variability)
from $-1$ to $0$ -- since that framing is not really climate-skeptical,
just non-climate -- leaves both findings essentially unchanged (year trend
$\beta=0.0102$, $p<0.001$; outlet lean $\beta=-0.0135$, $p<0.001$). Third,
dropping the 10 category-E (contested) articles entirely, rather than
scoring them $0$, again leaves both findings unchanged (year trend
$\beta=0.0130$, $p<0.001$; outlet lean $\beta=-0.0135$, $p=0.001$).

\paragraph{Sample restrictions.} Restricting the disaster-level sample to
the 138 disasters with $N_d \geq 3$ (so the index is not a single
article's score) leaves the year trend significant
($\beta=0.0112$, $p<0.001$). Restricting the outlet-lean test to
post-2010 articles only -- since the pre-2010 era contributes just 37 of
209 covered disasters and could otherwise be driving the result --
likewise leaves it significant ($\beta=-0.0149$, $p=0.001$, $N=779$).

\paragraph{Placebo checks.} If outlet political lean is capturing
climate-framing specifically, it should not also predict things that have
nothing to do with framing. Two such checks are run. First, lean against
article word count: this is \emph{not} a clean null
($\beta=-18.7$ words per Bias-score point, $p=0.075$) -- more
right-leaning outlets in this sample write modestly shorter articles, so
part of the lean-attribution relationship plausibly reflects a generic
length/depth difference across outlets rather than climate-framing
decisions specifically, a limitation we note rather than dismiss. Second,
lean against the covered disaster's (NOAA-reported, not outlet-chosen)
log-cost and log-deaths: both are clean nulls ($p=0.685$ and $p=0.300$
respectively), meaning more right- or left-leaning outlets are not
systematically selecting into covering more or less severe disasters --
the lean effect is not a selection-into-coverage artifact, even if it may
partly reflect a style/depth difference per the word-count check.

\paragraph{National-only vs.\ local-only subsamples.} A different question
from Section~\ref{sec:outlet_results} (does national status itself predict
the score) is whether the two headline effects are \emph{homogeneous}
across national and local coverage, or concentrated in one segment.
Re-estimating each headline regression separately by outlet type: the year
trend holds in both ($\beta=0.0174$, $p<0.001$ among national-outlet
articles; $\beta=0.0107$, $p=0.007$ among local-outlet articles). The
outlet-lean effect does not generalize the same way -- it is strongly
significant among national-outlet articles ($\beta=-0.0308$, $p=0.001$,
$N=687$) but statistically indistinguishable from zero among local-outlet
articles ($\beta=-0.0025$, $p=0.69$, $N=202$). This qualifies the
political-lean finding meaningfully: it is robust to every measurement and
specification choice tried above, including the demanding disaster-
fixed-effects test, but appears to be concentrated in (or driven by)
national/wire-service coverage rather than holding uniformly across local
outlets as well -- though the local subsample is also a third the size of
the national one, so this could partly reflect lower power rather than a
true zero effect locally. We report this as an open qualification rather
than resolve it, since distinguishing the two would require a larger
sample of lean-rated local outlets than is currently available.

\begin{sloppypar}
\paragraph{Ordered logit/probit (functional-form check).} Every result
above treats $a_{i,d}\in\{-1,0,0.5,1\}$ as cardinal via OLS -- a real
assumption, since the gap between ``no claim'' (C) and ``hedged mention''
(B) need not be the same size as the gap between B and ``explicit
attribution'' (A). Dropping this assumption entirely and re-estimating
both headline relationships as an ordered outcome (collapsing C and E to
one level, since both score $0$: $D < \{C,E\} < B < A$), via both ordered
logit and ordered probit, leaves both findings significant under all four
combinations: the year trend (logit $\beta=0.104$, $p<0.001$; probit
$\beta=0.055$, $p<0.001$) and the outlet-lean effect (logit
$\beta=-0.085$, $p=0.004$; probit $\beta=-0.043$, $p=0.006$). Neither
result is an artifact of the cardinal-scaling assumption.\footnote{The
ordered-model standard errors here are not clustered by disaster (the
\texttt{statsmodels} ordered-model implementation does not support
cluster-robust SE), so these $p$-values are an approximate check on sign
and significance rather than directly comparable to the clustered-SE
results elsewhere in this section.}
\end{sloppypar}

\paragraph{Year $\times$ outlet-lean interaction.} Whether the
political-lean gap in framing has widened or narrowed over the sample
period is a separate question from its average magnitude. Adding an
$\text{OutletLean}_i \times \text{Year}_d$ interaction term to the
article-level model, the interaction is not statistically significant
either alone ($\beta=-0.00028$, $p=0.61$) or with the full set of controls
($\beta=-0.00030$, $p=0.66$). The point estimate implies a slightly more
negative (i.e.\ modestly larger-magnitude) lean effect by the end of the
sample than the beginning, but this is not distinguishable from no change
given the standard error -- we find no evidence the ideological gap in
attribution framing has systematically widened or narrowed over
2000--2024.

Appendix Table~\ref{tab:robustness_summary} summarizes all checks. With
the exception of the word-count placebo and the local-only outlet-lean
subsample, every alternative specification and sample restriction leaves
both headline results intact.

\paragraph{Audit results.} All results in Section~\ref{sec:results} are
estimated from the full-sample LLM coding pass. The independent post-hoc
audit sample (100 articles, described in Section~\ref{sec:methods}) was
hand-coded and compared against the LLM's codes. Relevance
agreement was 99\% raw / $\kappa = 0.96$ (target $\kappa \geq 0.7$);
Attribution agreement, on the 86 articles both raters independently
classified as Relevant, was 95.3\% raw, unweighted $\kappa = 0.93$,
linearly weighted $\kappa_w = 0.96$, and quadratically weighted
$\kappa_w = 0.98$ (targets $\kappa \geq 0.7$, ideally $\kappa_w \geq 0.8$),
with a mean absolute score difference of $0.023$ on the $[-1,1]$ scale.
All four Attribution disagreements were between adjacent categories (two
A/B, one B/C, one B/E); none spanned a larger gap. These confirm, on an
independent sample, the high reliability already suggested by the
calibration-sample agreement ($\kappa=0.81$ for Relevance, $\kappa=1.00$
for Attribution) -- the results in Section~\ref{sec:results} can be read
as confirmed rather than provisional.


% ===================================================================
\section{Discussion}
% ===================================================================

\subsection*{The Year Trend and the 2010 Structural Break}

The most robust finding in this paper is a steady upward trend in
anthropogenic attribution framing: covered disasters average an
Attribution Index of $-0.105$ before 2010 and $+0.057$ from 2010
onward, a shift of more than 0.16 index points that persists under
every robustness check we apply (Appendix
Table~\ref{tab:robustness_summary}).
This timing aligns closely with the structural break identified by
\citet{chinn2020}, who find using unsupervised text scaling (Wordfish)
that partisan \emph{polarization} in U.S.\ climate news language was
relatively flat through the 2000s before undergoing a discrete jump in
2010--2011.  Our result complements theirs: not only did the language
of climate coverage polarize around 2010, but the specific practice of
linking extreme weather events to anthropogenic causes --- which had
been net-negative through the early 2000s --- shifted decisively toward
net-positive after that point.  The negative pre-2010 mean in our data
is also consistent with the ``balance as bias'' era documented by
\citet{boykoff2004}: their finding that the U.S.\ prestige press over
1988--2002 systematically gave equal weight to the scientific consensus
and a small minority of skeptics --- producing net-neutral or
net-skeptical aggregate framing --- aligns with our pre-2010 mean
Attribution Index of $-0.105$.  \citet{boykoff2007} provides the
temporal bridge: extending the analysis through 2006, he documents that
the balance norm eroded sharply in U.S. prestige-press coverage by
2005--2006, with balanced coverage falling from 36.6\% in 2003
(statistically significant, $p<0.001$) to 3.3\% in 2006 (no longer
significant) --- a convergence toward scientific consensus that preceded
our post-2010 shift to net-positive attribution framing by roughly
five years, and that together with \citet{chinn2020}'s structural break
in 2010--2011 traces a coherent timeline from the balance-norm era to
the present.

This paper identifies the \emph{when} of the post-2010 shift more
cleanly than the \emph{why}; several mechanisms are plausible.
First, the scientific literature on extreme-event attribution expanded
substantially after roughly 2011, when formal attribution science
(asking whether anthropogenic forcing changed the probability of a
specific event) began receiving mainstream scientific attention; growing
scientific output on individual events plausibly gave journalists more
attributable claims to report.  Second, \citet{chinn2020} document that
political actors outnumber scientists as quoted sources in U.S.\ climate
news from 2006 onward, suggesting that supply-side editorial choices ---
which voices are treated as authoritative --- shift the content of
coverage independently of the underlying science.  Third, the period
after 2009--2010 saw a series of individually prominent extreme events
(major flooding, record droughts, Hurricane Sandy) that attracted
sustained national attribution commentary, potentially establishing a
template for how subsequent events were covered.  Disentangling these
mechanisms would require event-level attribution-science publication
dates matched to coverage timing, which is beyond the scope of the
present paper but is a natural extension.

\subsection*{Outlet Political Lean: Supply, Demand, and the
National--Local Asymmetry}

The second main finding is that more right-leaning outlets (by Ad
Fontes Bias score) produce systematically lower attribution scores,
with an estimated gradient of $-0.011$ per Bias-score point
($p=0.013$, clustered SE), robust to disaster fixed effects, both lean
measures, all attribution-measure variants, and the ordered-logit
functional form.  Two broad explanations are consistent with this
pattern: a supply-side story in which editorial decisions reflect the
ideological preferences of editors, owners, or reporters; and a
demand-side story in which outlets cater to the beliefs of their
readers.

\citet{gentzkow2010} provide relevant evidence on this question in the
context of general media slant: using the frequency of partisan phrases
in congressional speech to construct outlet slant indices, they find
that approximately 20\% of cross-outlet variation in slant is explained
by reader demographic demand, while owner identity explains
substantially less.  If the same demand-side logic applies to
attribution framing specifically, much of the lean gradient we document
may reflect outlets supplying the causal narratives their audiences
expect: right-leaning outlets serve audiences who, consistent with
\citet{carmichael2017}'s evidence on partisan media and climate concern
and \citet{mccright2011}'s documentation of growing partisan
polarization in climate beliefs over 2001--2010 (ideological gap
widening from 18 to 44 percentage points over the decade), are less
receptive to anthropogenic attribution framing.

Content-analytic evidence from the same newspapers we study lends
additional support to this outlet-lean gradient.  \citet{feldman2017}
find, for 2006--2011, that the \emph{Wall Street Journal} --- the
prestige daily with the most right-leaning editorial orientation ---
was significantly less likely than the \emph{New York Times},
\emph{Washington Post}, and \emph{USA Today} to discuss climate
impacts, least likely to frame climate change as an immediate or
domestic threat, and most likely to use conflict and negative-economic
frames for climate actions.  Our outlet-lean result extends this
single-paper comparison: the lean-attribution gradient holds
continuously across a wide range of outlets (from wire services to
regional papers), persists over a substantially longer sample period
(2000--2024), and survives controls for disaster type, year, and
disaster fixed effects.

The national--local asymmetry qualifies this picture importantly.  The
lean gradient is strongly significant among national-outlet articles
($\hat\beta = -0.031$, $p=0.001$, $N=687$) but is statistically
indistinguishable from zero among local-outlet articles
($\hat\beta = -0.003$, $p=0.69$, $N=202$).  The local subsample is
about one-third the size of the national one, so part of this
difference may reflect lower power rather than a true zero effect among
local outlets.  Nonetheless, it raises the possibility that local
papers --- which serve geographically specific audiences who may have
directly experienced an event --- behave differently from national and
wire outlets when deciding how to frame causation.  \citet{ellger2024}
show, in the context of German newspaper exits, that local papers play
a distinct informational role that is not substitutable by national
coverage: when a local paper exits, consumers shift to national
tabloids that are more partisan and less locally focused, and political
polarization rises as a consequence.  A symmetric reading of their
result applied here would suggest that local papers may moderate
partisan attribution framing relative to national outlets, and that the
ongoing contraction of local print journalism could shift the
aggregate information environment further toward the nationally
partisan lean gradient we document.

The non-significance of the year $\times$ lean interaction
($p = 0.61$ in the simple model, $p = 0.66$ with full controls) adds a
further nuance: the ideological gap in attribution framing has not
significantly widened or narrowed over the sample period.  Whatever
mechanism generates the lean gradient appears to have been stable
across the 2000--2024 window, even as the overall level of attribution
framing rose.

\subsection*{Null Results on Disaster Characteristics}

A striking feature of the results is the near-complete absence of
predictive power from the disaster's own physical characteristics.
Disaster type, CPI-adjusted cost, death toll, and duration are all
statistically indistinguishable from zero in the disaster-level
specification (Table~\ref{tab:main_results}), and a one-way ANOVA
across the seven NOAA event types confirms no significant variation
($F = 0.32$, $p = 0.93$; Appendix Table~\ref{tab:type_means}).  This null
result on disaster type is consistent with \citet{boudet2020}, who
find that even events with a high scientific attribution confidence
score (such as extreme heat) do not reliably generate more
attribution discussion at the local level than events with lower
attribution confidence.  In both their study and ours, the content of
coverage appears more responsive to contextual factors outside the
event itself --- timing, outlet identity, and the broader information
environment --- than to the physical or scientific characteristics of
the disaster being covered.

This finding also has a methodological implication.  If disaster type
and severity do not predict attribution framing in the cross-section,
it is unlikely that the composition of NOAA events across years (e.g.,
more tropical cyclones in some years than others) is what drives the
year trend.  The year trend appears to be a genuine temporal shift in
how coverage is framed, not a composition artifact.

\subsection*{Coverage Volume}

Coverage volume ($N_d$, the number of Relevant articles per disaster)
enters the disaster-level regression positively and significantly
($\hat\beta = 0.009$, $p = 0.036$): disasters that attract more press
coverage skew modestly more anthropogenic on average.  This
association admits two interpretations that are difficult to
distinguish in observational data.  On the first, a high climate-angle
salience increases both the volume of coverage and the probability of
attribution framing --- the same editorial judgment that generates more
articles also generates more attribution language, so the two are joint
outcomes of latent salience.  On the second, with more articles
covering a given event there is a higher probability that at least one
of them introduces attribution framing that then pulls the disaster-
level average above zero, a purely mechanical relationship.  Given this
ambiguity, we control for $N_d$ throughout the analysis but do not
interpret the coefficient causally.

\subsection*{Implications for Risk Perception and Policy}

The pattern we document --- near-zero average Attribution Index, a
rising trend post-2010, and a partisan lean gradient concentrated in
national coverage --- has several implications for the larger research
programs on disaster risk perception and preparedness.

This complements \citet{gallagher2014}'s finding that the salience
channel from flood coverage to insurance take-up operates at the same
magnitude regardless of the disaster's dollar cost.  If attribution
framing further shapes how risk is perceived -- whether a disaster
reads as a precedent for future events or as a one-off -- then the
near-zero and historically non-anthropogenic framing we document for
most of the sample period may have dampened the longer-run risk
signaling that coverage could otherwise have provided.

This complements \citet{magontier2020}'s finding that storm coverage
\emph{volume} drives FEMA hazard-mitigation investment: whether coverage
\emph{content} --- natural-variability framing versus anthropogenic
framing --- independently modulates that preparedness response remains
untested, since it would require the mitigation-project data to be
crossed with article-level attribution codes at the event level.

At the individual level, \citet{egan2012} document that personal
exposure to anomalously high local temperatures increases the
probability that a respondent believes there is solid evidence of
warming, with the effect largest among those who are less educated and
who have weak attachments to political parties.  The fact that this
effect decays within 12 days suggests that individual-level risk
updating from direct experience is short-lived, which raises the
relative importance of media framing as a more durable information
source.  If right-leaning coverage is systematically less likely to
frame extreme events as anthropogenically caused --- as we find --- then
the populations that rely most on national right-leaning outlets for
information about climate may receive a systematically different signal
about the causes of the disasters they experience.

\subsection*{Limitations}

Several limitations should inform how the findings are read.  First,
the Attribution Index is an average over a small and potentially
unrepresentative sample of articles: 33.2\% of NOAA events have no
Relevant coverage at all, and the articles we do have are concentrated
in large national and wire-service outlets.  Whether the pattern we
document in covered disasters generalizes to the median NOAA event is
uncertain.

Second, the outlet-lean gradient is identified from the 33 domains with
Ad Fontes or AllSides ratings, covering 71\% of Relevant articles but
concentrated in the largest outlets.  The local-outlet subsample with
lean ratings ($N=202$ articles) is small enough that the null effect
there could reflect limited power rather than a true zero.

Third, the LLM coding procedure, while validated by the post-hoc audit
($\kappa = 0.93$ unweighted, $\kappa_w = 0.96$--$0.98$ weighted),
relies on a single model and a single codebook.  The four
Attribution disagreements in the audit were all on adjacent-category
boundaries (A/B, B/C, B/E), which is the kind of fuzziness inherent
in any ordinal text classification scheme, but it implies that the
index is somewhat less reliable near the B/C boundary than at the
extremes.

Fourth, the outlet-lean finding, while robust across all specification
checks, is correlational.  The most we can say is that more
right-leaning outlets produce less anthropogenic attribution framing
within the same disaster and year; we cannot rule out that a common
cause drives both outlet lean and coverage decisions without controlling
for it.  The word-count placebo ($p = 0.075$) raises the additional
possibility that some of the lean gradient reflects a generic
length-and-depth difference across outlet types rather than a
climate-framing decision specifically.


% ===================================================================
\section{Conclusion}
% ===================================================================

No prior study has systematically measured how U.S.\ newspapers frame
the causal origins of billion-dollar climate disasters at scale; this
paper does so.  Using an LLM-based coding pipeline validated against an
independent human audit ($\kappa = 0.96$ for Relevance,
$\kappa_w = 0.96$--$0.98$ for Attribution), we assign each of 1,252
Relevant articles --- collected across all 313 NOAA billion-dollar
disaster events from 2000 to 2024 --- to one of five attribution
categories ranging from explicit anthropogenic framing ($+1$) to
explicit non-climate framing ($-1$), and aggregate these to a
disaster-level Attribution Index.

Three findings stand out.  First, silence is the modal outcome: 74.8\%
of Relevant articles contain no causal claim about a disaster's origin
(category C), leaving the mean disaster-level Attribution Index at just
0.028 --- statistically distinguishable from zero but economically
close to it.  This is not a media environment split between climate
advocates and skeptics; it is one in which most disaster coverage
simply avoids the question of cause.

Second, the trend runs in one direction: coverage has grown steadily
more willing to name a cause, adding roughly 0.013 index points a year
($p < 0.001$) and moving from a net-negative mean before 2010 to a
net-positive one after --- a shift that lines up with the 2010--2011
structural break \citet{chinn2020} document in partisan climate
discourse, and that survives every functional-form and
sample-restriction check we run.  Even by the end of the sample, though,
pure event reporting remains the dominant category.

Third, who publishes the coverage matters: more right-leaning outlets
(by Ad Fontes Bias score) systematically produce lower Attribution
Index values ($\hat\beta = -0.011$ per Bias-score point, $p = 0.013$,
clustered SE).  The gradient holds up against disaster fixed effects,
both lean-rating services, every attribution-measure variant, and an
ordered-logit functional-form check, but it is concentrated in national
and wire-service coverage and is not statistically detectable among
local outlets, where the subsample is small enough that limited power
cannot be ruled out.  The gap itself has neither widened nor narrowed
over the sample period: the year $\times$ lean interaction is not
significant, so whatever generates the ideological gradient appears to
have held steady even as overall attribution framing rose.

None of this traces back to what actually happened on the ground:
disaster type, CPI-adjusted cost, death toll, duration, coastal
geography, and the political lean of affected states are all
uninformative about how coverage frames causation.  That null result
matters methodologically --- it rules out composition as the driver of
the patterns above: the mix of events did not shift after 2010 in a way
that would mechanically raise attribution framing, and right-leaning
outlets are not simply covering less climate-salient disasters.

\paragraph{Contributions.}  The paper makes three contributions.
Methodologically, it demonstrates that a carefully validated LLM coding
pipeline --- following the three-step text-as-data framework of
\citet{gentzkow2019} --- can reliably classify fine-grained causal
attribution in news text at a scale ($N > 1{,}000$ articles) that
would be prohibitively costly to hand-code, with inter-rater reliability
exceeding the targets set in advance.  Substantively, it produces the
first panel of attribution framing across the full universe of
NOAA billion-dollar disasters, enabling the year-trend and outlet-lean
analyses that are the paper's headline results.  Empirically, it
connects the framing literature to the economics literature on disaster
salience: given that \citet{gallagher2014} shows media coverage of
flood events materially shifts risk-protective behavior even in
non-affected communities in the same media market, the attribution
content of that coverage is a potentially important and previously
unmeasured input into how households and governments update their
beliefs and preparedness decisions.

\paragraph{Limitations and future directions.}  Several limitations
point toward natural extensions.  The sample is constrained to events
meeting NOAA's billion-dollar threshold (in real 2024 dollars),
excluding the much larger universe of smaller disasters; whether
attribution framing patterns differ for sub-threshold events is
unknown.  The outlet universe, while broad (899 verified domains), is
biased toward outlets large enough to be indexed by Google's search
engine; smaller community papers are underrepresented, and the lean-
rated subsample is further concentrated in the largest national outlets.
Extending the analysis to a more representative sample of local
coverage --- particularly as the local-newspaper industry continues to
contract \citep{ellger2024} --- would clarify whether the
national--local asymmetry in lean effects we document is a durable
structural feature or a power artifact.  A related extension is to
expand the article universe itself through a more comprehensive
outlet-discovery and retrieval pipeline.  The current sample is
constrained by the coverage of the search infrastructure used to
collect candidate articles, which skews toward outlets with high web
presence and may underrepresent regional and community newspapers that
are indexed less systematically.  Crucially, the validated LLM coding
procedure developed here --- with inter-rater reliability established
at $\kappa = 0.96$ for relevance and $\kappa_w = 0.96$--$0.98$ for
attribution --- would transfer directly to a larger corpus without
modification.  The pipeline's scalability means that improvements
upstream in outlet discovery and article retrieval directly translate
into a richer panel of attribution framing data, without requiring any
re-validation of the coding step itself.

Future work could also link the Attribution Index directly to
downstream outcomes.  \citet{magontier2020} shows that storm media
coverage volume predicts FEMA mitigation investment; a natural next
step is to ask whether the \emph{content} of coverage --- anthropogenic
versus natural-variability framing --- independently predicts
preparedness, insurance take-up, or political responses.  On the
methodology side, an obvious extension is to apply the same pipeline
to other countries' media markets, where the mapping between outlet
lean and attribution framing may differ from the U.S.\ pattern given
different partisan polarization dynamics and media market structures.

Finally, the year-trend finding raises a question the paper cannot
answer: how much of the post-2010 rise in attribution framing reflects
the growth of formal event-attribution science (which gave journalists
attributable claims to cite), how much reflects shifting elite cues
from politicians and advocacy organizations, and how much reflects
demand-side changes in audience interest.  Answering this would require
matching attribution-science publication dates, politician attribution
statements, and audience engagement metrics to article-level coverage
timing --- a demanding but tractable research design that would
substantially clarify the mechanisms behind the trend documented here.


% ===================================================================
\newpage
\bibliographystyle{plainnat}
\bibliography{references}


% ===================================================================
\appendix
\section{Data Appendix}
% ===================================================================

\subsection{Data and Code Availability}
\label{sec:data_availability}

\begin{sloppypar}
All data and code underlying this paper are available in the project's
GitHub repository,
\href{https://github.com/Hariaksha/climate-disaster}{github.com/Hariaksha/climate-disaster}.
The central data file -- the full \texttt{Master} sheet of article-level
codes, the \texttt{All Disasters} disaster-level table, the
\texttt{Calibration} and \texttt{Audit Sample} sheets, and the
\texttt{Domain Ratings} outlet-lean table -- is
\href{https://github.com/Hariaksha/climate-disaster/blob/main/attribution.xlsx}{\texttt{attribution.xlsx}}.
The Stage 3 regression analysis reported in
Sections~\ref{sec:results}--\ref{sec:results_robustness} is fully
reproducible from
\href{https://github.com/Hariaksha/climate-disaster/blob/main/code/Stage3_Regressions.ipynb}{\texttt{code/Stage3\_Regressions.ipynb}},
and the attribution-coding codebook is
\href{https://github.com/Hariaksha/climate-disaster/blob/main/paper/attribution_codebook.md}{\texttt{paper/attribution\_codebook.md}}.
Note that \texttt{attribution.xlsx} includes full retrieved article text
(\texttt{Master} column \texttt{Full Article Text}), which is copyrighted
newspaper content; this is shared here for replication purposes, but
re-publication of that column specifically beyond replication use is not
authorized by the original publishers.
\end{sloppypar}

\subsection{Disaster Event List}

The full list of 313 NOAA events, with assigned states/regions and PSE
engine coverage, is the \texttt{All Disasters} sheet of
\href{https://github.com/Hariaksha/climate-disaster/blob/main/attribution.xlsx}{\texttt{attribution.xlsx}}
(not reproduced in full here; see Section~\ref{sec:data} for summary
statistics).

\subsection{Newspaper Source List}

\begin{sloppypar}
The 19 Programmable Search Engines and their domain coverage are
documented in
\href{https://github.com/Hariaksha/climate-disaster/blob/main/data/Access\%20World\%20News\%20Database.xlsx}{\texttt{data/Access World News Database.xlsx}};
the resulting outlet-level political-lean ratings (AllSides, Ad Fontes)
used in Section~\ref{sec:outlet_results} are the \texttt{Domain Ratings}
sheet of
\href{https://github.com/Hariaksha/climate-disaster/blob/main/attribution.xlsx}{\texttt{attribution.xlsx}}.
\end{sloppypar}

\subsection{Additional Figures}

The two figures below supplement the main-text discussion of the
Attribution Index's cross-sectional distribution (Section~\ref{sec:results},
Descriptive Patterns) and of the outlet political-lean result
(Section~\ref{sec:outlet_results}, Variation by Outlet Characteristics).

\begin{figure}[ht]
\centering
\includegraphics[width=0.7\textwidth]{figures/attribution_index_hist.pdf}
\caption{Distribution of the Disaster-Level Attribution Index ($N=209$
covered disasters). Vertical line at the mean ($0.028$); dotted line at
zero. Discussed in Section~\ref{sec:results} alongside the year-trend
result in Figure~\ref{fig:trend_over_time}.}
\label{fig:index_hist}
\end{figure}

\begin{figure}[ht]
\centering
\includegraphics[width=0.75\textwidth]{figures/outlet_lean_scatter.pdf}
\caption{Outlet Political Lean vs.\ Article Attribution Score ($N=889$
articles from individually-rated domains). Points are jittered on the
$y$-axis for visibility, since the underlying score only takes four
values ($-1, 0, 0.5, 1$); diamonds are means within 8 equal-width bins of
Ad Fontes Bias score. Discussed in Section~\ref{sec:outlet_results}.}
\label{fig:outlet_lean_scatter}
\end{figure}

\FloatBarrier
\subsection{Additional Tables}

The four tables below supplement the main-text discussion of variation by
disaster type (Section~\ref{sec:results}), the outlet political-lean
result (Section~\ref{sec:outlet_results}), and the full robustness battery
(Section~\ref{sec:results_robustness}).

\begin{table}[ht]
\centering
\caption{Mean Attribution Index by Disaster Type}
\label{tab:type_means}
\begin{threeparttable}
\begin{tabular}{lcccc}
\toprule
Disaster Type & $N$ & Mean & Median & SD \\
\midrule
Severe Storm      & 89 & 0.005 & 0.000 & 0.236 \\
Tropical Cyclone  & 45 & 0.055 & 0.022 & 0.113 \\
Flooding          & 25 & 0.028 & 0.000 & 0.175 \\
Drought           & 18 & 0.026 & 0.017 & 0.476 \\
Wildfire          & 18 & 0.048 & 0.021 & 0.315 \\
Winter Storm      & 12 & 0.074 & 0.000 & 0.148 \\
Freeze            &  2 & 0.000 & 0.000 & 0.000 \\
\bottomrule
\end{tabular}
\begin{tablenotes}
\small
\item Notes: One-way ANOVA across the seven types: $F=0.32$, $p=0.93$.
\end{tablenotes}
\end{threeparttable}
\end{table}

\begin{table}[ht]
\centering
\caption{Robustness Check Summary}
\label{tab:robustness_summary}
\begin{threeparttable}
\begin{tabular}{lcc}
\toprule
Check & Estimate & $p$-value \\
\midrule
\multicolumn{3}{l}{\emph{Year trend}} \\
\quad Baseline                              & $0.0129$/yr & $<0.001$ \\
\quad Binary attribution measure            & $0.0100$    & $<0.001$ \\
\quad D-subtype recoded                     & $0.0102$    & $<0.001$ \\
\quad Contested (E) excluded                & $0.0130$    & $<0.001$ \\
\quad $N_d \geq 3$ only                     & $0.0112$    & $<0.001$ \\
\quad National-outlet articles only         & $0.0174$    & $<0.001$ \\
\quad Local-outlet articles only            & $0.0107$    & $0.007$ \\
\quad Ordered logit                         & $0.1044$    & $<0.001$ \\
\quad Ordered probit                        & $0.0546$    & $<0.001$ \\
\midrule
\multicolumn{3}{l}{\emph{Outlet political lean (Ad Fontes Bias score)}} \\
\quad Baseline (multivariate, clustered SE) & $-0.0112$   & $0.013$ \\
\quad Bivariate (no controls)               & $-0.0132$   & $0.002$ \\
\quad Disaster fixed effects                & $-0.0115$   & $0.046$ \\
\quad Binary attribution measure            & $-0.0177$   & $<0.001$ \\
\quad D-subtype recoded                     & $-0.0135$   & $0.001$ \\
\quad Contested (E) excluded                & $-0.0135$   & $0.001$ \\
\quad Post-2010 only                        & $-0.0149$   & $0.001$ \\
\quad National-outlet articles only         & $-0.0308$   & $0.001$ \\
\quad Local-outlet articles only            & $-0.0025$   & $0.689$ \\
\quad Ordered logit                         & $-0.0850$   & $0.004$ \\
\quad Ordered probit                        & $-0.0434$   & $0.006$ \\
\quad $\times$ Year interaction             & $-0.00028$  & $0.611$ \\
\midrule
\multicolumn{3}{l}{\emph{Placebo checks (outlet lean predicting non-framing outcomes)}} \\
\quad Word count                            & $-18.69$    & $0.075$ \\
\quad log(Cost) of covered disaster          & $-0.012$    & $0.685$ \\
\quad log(Deaths+1) of covered disaster      & $-0.034$    & $0.300$ \\
\bottomrule
\end{tabular}
\begin{tablenotes}
\small
\item Notes: Year-trend checks are disaster-level OLS, robust (HC1) SE,
except the ordered logit/probit rows, which are article-level (ordinal
attribution category regressed on year, non-clustered MLE SE). Outlet-
lean and placebo checks are article-level OLS, SE clustered by disaster
(or within-disaster-demeaned with no intercept, for the fixed-effects
row), except the ordered logit/probit and interaction rows (non-clustered
MLE SE and clustered OLS, respectively -- see notebook for exact
specifications). Full detail and code for every row in
\texttt{code/Stage3\_Regressions.ipynb}, Section K.
\end{tablenotes}
\end{threeparttable}
\end{table}

\begin{table}[ht]
\centering
\caption{Full Disaster-Type Coefficients, Disaster-Level Model (Column 3 of Table~\ref{tab:main_results})}
\label{tab:appendix_type_full}
\begin{threeparttable}
\begin{tabular}{lcc}
\toprule
Disaster Type (ref.\ = Severe Storm) & Coef. & SE \\
\midrule
Drought          & 0.0050  & (0.1260) \\
Flooding         & 0.0195  & (0.0488) \\
Freeze           & 0.1119  & (0.0667) \\
Tropical Cyclone & 0.0106  & (0.0556) \\
Wildfire         & 0.0407  & (0.0903) \\
Winter Storm     & 0.0790  & (0.0695) \\
\bottomrule
\end{tabular}
\begin{tablenotes}
\small
\item Notes: Robust (HC1) SE. None of the six type dummies is individually
significant at $p<0.10$. Companion coefficients (cost, deaths, duration,
geography, year, $N_d$) are in Table~\ref{tab:main_results}, column 3.
\end{tablenotes}
\end{threeparttable}
\end{table}

\begin{table}[ht]
\centering
\caption{Full Disaster-Type Coefficients, Article-Level Model (Column 2 of Table~\ref{tab:outlet_results})}
\label{tab:appendix_article_full}
\begin{threeparttable}
\begin{tabular}{lcc}
\toprule
Disaster Type (ref.\ = Severe Storm) & Coef. & SE \\
\midrule
Drought          & 0.3412$^{***}$ & (0.1054) \\
Flooding         & 0.1020$^{**}$  & (0.0518) \\
Freeze           & 0.1442$^{**}$  & (0.0708) \\
Tropical Cyclone & 0.1203$^{**}$  & (0.0478) \\
Wildfire         & 0.2411$^{***}$ & (0.0776) \\
Winter Storm     & 0.0663         & (0.0481) \\
\bottomrule
\end{tabular}
\begin{tablenotes}
\small
\item Notes: $^{*}p<0.10$, $^{**}p<0.05$, $^{***}p<0.01$. SE clustered by
disaster. Unlike the disaster-level model, disaster type is strongly
significant here: at the article level, with disaster-level confounds
(cost, geography, etc.) absorbed differently and $N$ roughly 4--6$\times$
larger, droughts, wildfires, tropical cyclones, floods, and freezes all
score significantly higher than severe storms. Companion coefficients
(National, Ad Fontes Bias, Year) are in Table~\ref{tab:outlet_results},
column 2. The contrast with Table~\ref{tab:appendix_type_full} (where type
is insignificant) reflects the different unit of analysis and sample
restriction (889 articles from rated domains vs.\ 201 disasters) and
should not be over-read as a contradiction.
\end{tablenotes}
\end{threeparttable}
\end{table}

The full Stage 3 analysis -- dataset construction, all regressions and
cross-tabs reported in Section~\ref{sec:results}, plus the robustness
checks and placebo tests in Section~\ref{sec:results_robustness} -- is
reproducible from
\href{https://github.com/Hariaksha/climate-disaster/blob/main/code/Stage3_Regressions.ipynb}{\texttt{code/Stage3\_Regressions.ipynb}}
(see Section~\ref{sec:data_availability} for the full data/code
availability statement).

\end{document}
